% !TEX root = ../../thesis.tex

\section{Introduction}

The construction of minimal hypersurfaces via min-max methods involves two essentially independent components: a \emph{sweepout framework} that produces a candidate varifold, and a \emph{regularity theory} that promotes this varifold to a smooth hypersurface.

In the classical Almgren--Pitts theory~\cite{Alm65, Pit81}, the sweepout framework operates in the space of integral currents: a combinatorial selection procedure extracts a min-max sequence whose limit satisfies the \emph{almost minimizing} property---the varifold is almost area-minimizing in small annuli. This property feeds directly into the Schoen--Simon regularity theory~\cite{SS81}, which requires either the almost minimizing condition or, more generally, a strong a priori bound on the singular set ($\mathcal{H}^{n-2}(\sing V) < \infty$). Because the construction uses integral currents, integrality of the limit is built in from the outset.

Following the ideas of Simon--Smith~\cite{Smi83}, Colding--de~Lellis~\cite{CL03} and de~Lellis--Tasnady~\cite{DLT13} replaced the discrete setting of Almgren--Pitts with continuous sweepouts by Caccioppoli sets, considerably simplifying the combinatorial machinery. However, they still rely on the same regularity pipeline: an almost minimizing property plus Schoen--Simon.

Chodosh--Liokumovich--Spolaor~\cite{CLS22} introduced a fundamentally different sweepout framework based on \emph{non-excessive sweepouts}---optimal nested sweepouts whose critical slices cannot be locally replaced by competitors of strictly smaller mass. This framework is considerably more streamlined: it works directly with Caccioppoli sets, avoids the combinatorial machinery of Almgren--Pitts, and replaces the almost minimizing property with a more intuitive geometric condition---one-sided homotopic minimization. However, the regularity step in~\cite{CLS22} still invokes the Almgren--Pitts pull-tight and regularity theory as a black box~\cite[Section~3]{CLS22}, leaving the non-excessive framework dependent on the classical machinery it was designed to simplify.

In this chapter, we explore the extent to which this dependence can be removed. We construct an alternative regularity pipeline that replaces the Almgren--Pitts combinatorial selection and Schoen--Simon regularity: stationarity follows from the pull-tight argument of Colding--de~Lellis~\cite{CL03}; integrality from the de~Lellis--Tasnady criterion~\cite[Proposition~A.1]{DLT13} or, when mass cancellation occurs, from a sheet-counting argument exploiting nestedness (Proposition~\ref{prop:integrality-nested}); and regularity from Wickramasekera's theorem~\cite{Wic14} for stable codimension~$1$ integral varifolds, whose hypotheses are naturally satisfied in the non-excessive setting.

Wickramasekera's theorem applies to stable integral varifolds satisfying the \emph{$\alpha$-structural hypothesis} (Definition~\ref{defn:alpha-structural}), which roughly asks that no singular point has a neighbourhood in which three or more hypersurface sheets meet along a common boundary. It strictly subsumes the Schoen--Simon theory, which requires either the almost minimizing condition or an a priori bound $\mathcal{H}^{n-2}(\sing V) < \infty$ on the singular set. The proof of Wickramasekera's theorem is substantially more involved than that of Schoen--Simon, requiring a simultaneous induction on two regularity results (the Sheeting Theorem and the Minimum Distance Theorem), together with coarse and fine blow-up analysis and Simon's $L^2$-estimates.

The $\alpha$-structural hypothesis is well-suited to the non-excessive setting: the non-excessive condition directly rules out configurations where three or more sheets meet, and it ensures stability of the limit varifold---exactly the inputs Wickramasekera's theorem requires. The nested sweepout machinery developed in Part~1 of this thesis provides the optimal nested volume-parametrized sweepouts that serve as input to the construction (see Theorem~\ref{thm:non-excessive-existence}). Combining these ingredients, we obtain:

\begin{theorem}[Existence of minimal hypersurfaces via non-excessive min-max]\label{thm:non-excessive-main}
    Let $(M^{n+1}, g)$ be a closed Riemannian manifold with $2 \leq n \leq 6$. Then the non-excessive min-max construction, using nested sweepouts as input and Wickramasekera's regularity theorem for the regularity step, produces a smooth, closed, embedded minimal hypersurface in $M$.
\end{theorem}

The proof proceeds in three steps. First, we establish integrality of the limit varifold: in the absence of mass cancellation, the de~Lellis--Tasnady criterion~\cite[Proposition~A.1]{DLT13} applies directly (Theorem~\ref{thm:integrality-mult-one}); when mass cancellation occurs, we exploit nestedness of the sweepout to count sheets and deduce integer density via a graphical decomposition argument (Proposition~\ref{prop:integrality-nested}). Second, we verify the $\alpha$-structural hypothesis by analyzing tangent cones at potential singular points---in the no-mass-cancellation case, the Caccioppoli set structure constrains tangent cones to wedges and folds; in the mass-cancellation case, higher-order junctions are also possible but are ruled out by area-minimizing replacements. Third, the non-excessive structure ensures stability (via one-sided homotopic minimization), so Wickramasekera's theorem~\cite{Wic14} yields smoothness.

The remaining open question is \emph{multiplicity one}: whether the smooth limit has multiplicity~$1$, which is equivalent to the absence of mass cancellation (Section~\ref{sec:multiplicity-one}).

The chapter is organized as follows. Section~\ref{sec:p2-prelim} collects the basic definitions: Caccioppoli sets, varifolds, sweepouts, width, and the integrality criterion of de~Lellis--Tasnady. Section~\ref{sec:non-excessive-prelim} reviews the non-excessive sweepout framework of~\cite{CLS22}: non-excessive ONVP sweepouts, the pull-tight argument, and homotopic minimizers as the geometric mechanism through which non-excessiveness controls regularity. Section~\ref{sec:integrality} establishes integrality of the limit varifold in both the no-mass-cancellation and mass-cancellation cases. Section~\ref{sec:regularity} verifies stability and the $\alpha$-structural hypothesis, and completes the proof of Theorem~\ref{thm:non-excessive-main} via Wickramasekera's theorem. Finally, Section~\ref{sec:multiplicity-one} discusses the multiplicity-one question.
